\documentclass{article}
\usepackage{mathnotes}

\title{Partial Fraction Decomposition}
\description{Complete guide to decomposing rational expressions into simpler fractions, covering linear and quadratic factors, Heaviside cover-up method, and differentiation techniques with worked examples.}
\source{title={Ordinary Differential Equations}, author={Morris Tenenbaum and Harry Pollard}, type={book}, isbn={978-0-486-64940-5}, section={Lesson 26A}, published={1985}}

\begin{document}

\section{Partial Fraction Decomposition}
The goal of partial fraction composition is to be able to rewrite a rational expression as a sum of simpler fractions.

\subsection{Conditions for application}

Partial fraction decomposition works when we have a rational expression of the form:

\[ f(x) = \frac{P(x)}{Q(x)} \]

where both $P(x)$ and $Q(x)$ are polynomials and the degree of $P(x)$ is strictly less than the degree of $Q(x)$.

If the degree of $P(x)$ is greater than or equal to the degree of $Q(x)$, perform long division first and work with the remainder term.

\subsection{General Approach to Partial Fraction Decomposition}

This approach always works as long as the conditions above are satisfied, but there are other ways; see below.

\begin{enumerate}
\item First, factor the denominator as completely as possible.
\end{enumerate}

\begin{enumerate}
\item Handle \textbf{linear factors of} $Q(x)$
\end{enumerate}

Let $x - r$ be a linear factor of $Q(x)$. Suppose that $(x - r)^m$ is the highest power of $x - r$ that divides $Q(x)$.

This factor will translate to the sum of $m$ partial fractions, as follows:

\[ \frac{A_1}{x-r} + \frac{A_2}{(x-r)^2} + \cdots + \frac{A_m}{(x - r)^m}. \]

Repeat this for each distinct linear factor of $Q(x)$.

\begin{enumerate}
\item Handle \textbf{irreducible quadratic factors of} $Q(x)$
\end{enumerate}

An irreducible quadtric factor is a quadratic expression that cannot be factored into linear factors with real coefficients.

Let $x^2 + px + q$ be an irreducible quadratic factor of $Q(x)$. Suppose $(x^2 +px + q)^n$ is the highest power of this factor thad divides $Q(x)$.

This factor will translate to the sum of $n$ partial fractions, as follows:

\[ \frac{B_1 x + C_1}{x^2+px+q} + \frac{B_2 x + C_2}{(x^2+px+q)^2} + \cdots + \frac{B_n x + C_n}{(x^2+px+q)^n} \]

Repeat this for each distinct irreducible quadratic factor of $Q(x)$.

\begin{enumerate}
\item Set the original fraction $\frac{P(x)}{Q(x)}$ equal to the sum of all the partial fractions from steps 2 and 3.
\end{enumerate}

\begin{enumerate}
\item Clear the fractions by multiplying by $Q(x)$. Distribute. Arrange the terms in decreasing powers of $x$. Group coefficients of the same power of $x$ by factoring out that power of $x$.
\end{enumerate}

\begin{enumerate}
\item Equate the coefficients of corresponding powers of $x$ and solve the resulting equations for the undetermined coefficients.
\end{enumerate}

\begin{enumerate}
\item We can now write the equation in the form
\end{enumerate}

\[ \frac{P(x)}{Q(x)} = \frac{A_1}{x-r} + \frac{A_2}{(x-r)^2} + \cdots + \frac{A_m}{(x - r)^m} + \frac{B_1 x + C_1}{x^2+px+q} + \frac{B_2 x + C_2}{(x^2+px+q)^2} + \cdots + \frac{B_n x + C_n}{(x^2+px+q)^n} \]

\emph{note}: why can we only deal with linear and irreducible quadratics? Can we always factor higher degree polynomials into terms that are either linear or quadratic?

\subsubsection{Example}

\emph{jmh: my material}.

Express

\[ \tag{a} \frac{3(1+x)}{1-x^3} \]

as a sum of partial fractions.

Note that the degree of the numerator is 1 while the degree of the denominator is 3, so we satisfy the conditions for performing partial fraction decomposition and may proceed.

First, factor the denominator (recognizing that it's a sum of cubes):

\[ \tag{b} \frac{3(1+x)}{(1-x)(x^2+x+1)} \]

The denominator is now composed of one linear factor in $x$: $1+x$, and one irreducible quadratic factor in $x$:  $x^2+x+1$. Both are raised on the first power.

Write the partial fraction with an undetermined coefficient for the linear factor as:

\[ \tag{c} \frac{A}{1-x} \]

Write the partial fraction with undetermined coefficients for the irreducible quadratic factor as:

\[ \tag{d} \frac{Bx + C}{x^2+x+1} \]

Sum (c) and (d) and set it equal to (b):

\[ \tag{e} \frac{3(1+x)}{(1-x)(x^2+x+1)} = \frac{A}{1-x} + \frac{Bx + C}{x^2+x+1} \]

Clear the fractions by multiplying both sides by the denominator of the left hand side:

\[ \tag{f} 3(1+x) = A(x^2+x+1) + (Bx+C)(1-x) \]

Distribute

\[ \tag{g} 3 + 3x = Ax^2 + Ax + A + Bx - Bx^2 +C - Cx \]

Sort by descending powers of $x$:

\[ \tag{h} 3x + 3 = Ax^2 - Bx^2 + Ax + Bx - Cx + A + C \]

Group coefficients of like powers of $x$ by factoring out the common power of $x$:

\[ \tag{i} 3x + 3 = x^2(A-B) + x(A+B-C) + (A+C) \]

The equation will be an identity when the coefficients of respective powers of $x$ on both sides are equal. Set them up and solve for them.

\[ \tag{j} \text{Coefficients of}~x^2: 0 = A - B \]

\[ \text{Coefficients of}~x^1: 3 = (A+B-C) \]

\[ \text{Coefficients of}~x^0: 3 = (A+C) \]

Solving these simulataneously we get $A=2, B=2, C=1$

Substituting these values into (e) we get:

\[ \tag{k} \frac{3(1+x)}{(1-x)(x^2+x+1)} = \frac{2}{1-x} + \frac{2x + 1}{x^2+x+1} \]

which is the desired sum of partial fractions.

\subsection{The Heaviside Cover-up Method}

This method only applies when the factor of $Q(x)$ are:

\begin{itemize}
\item linear
\item distinct
\item raised to the first power
\end{itemize}

That is, they are of the form $(x-r_1),(x-r_2),\cdots,(x-r_n)$.

But when it does work, it's a bit faster than the \emph{proper} method above.

\begin{enumerate}
\item First, factor the denominator as completely as possible. We now have a function in the form:
\end{enumerate}

\[ \frac{P(x)}{Q(x)} = \frac{P(x)}{(x-r_1)(x-r_2)\cdots(x-r_n)} \]

\begin{enumerate}
\item For each of the $n$ factors, find the corresponding factor $A_n$ by setting $x$ to $r_n$ and \emph{covering} the factor containing $r_n$. By covering, we mean \emph{remove} it as it if weren't there. For example, for $n = 1$:
\end{enumerate}

\[ A_n = \frac{P(r_1)}{\cancel{(x-r_1)}(r_1-r_2)\cdots(r_1-r_n)} \]

\begin{enumerate}
\item Repeat $n$ times to find the $n$ different factors, then you can write the function in decomposed format:
\end{enumerate}

\[ \frac{P(x)}{Q(x)} = \frac{A_1}{x-r_1} + \frac{A_2}{x-r_2} + \cdots + \frac{A_n}{x-r_n} \]

\subsection{Differentiation}
Like the Heaviside cover up method, this method only works when the roots are distinct. This method comes from Morris and Tenenbaum's ODE book, lesson 26A.

A fraction of the form $\frac{1}{f(x)}$ will have a partial fraction expansion of the form:

\[ \tag{a} \frac{1}{f(x)} = \frac{A_1}{x - r_1} + \frac{A_2}{x - r_2} + \cdots + \frac{A_k}{x - r_k} + \cdots + \frac{A_n}{x - r_n } \]

If we multiply $(a)$ by $(x - r_k)$. Since $f(r_k) = 0$, there results:

\[ \tag{b} \frac{x - r_k}{f(x) - f(r_k)} = \frac{A_1(x - r_k)}{x - r_1} + \cdots + A_k + \cdots + \frac{A_n(x - r_n)}{x - r_n} \]

Let $x \longrightarrow r_k$. The left side of $(b)$ will approach $\frac{1}{f'(r_k)}$ and its right side will approach $A_k$. Hence,

\[ \tag{c} A_k = \frac{1}{f'(r_k)}, \quad k = 1,2,\cdots,n. \]

Substituting $(c)$ in $(a)$ we get:

\[ \tag{d} \frac{1}{f(x)} = \frac{1}{f'(r_1)(x - r_1)} + \frac{1}{f'(r_2)(x - r_2)} + \cdots + \frac{1}{f'(r_n)(x - r_n)}. \]

\subsection{Other Methods}
Thomas and Finney cover a couple of other methods of determining the coefficients:

\begin{itemize}
\item Differentiating
\item Assigning small values to $x$ such as $x=0,\pm1,\pm2$ to get equations in $A$, $B$, and $C$.
\end{itemize}

\subsection{References}
\href{https://tutorial.math.lamar.edu/classes/calcii/partialfractions.aspx}{Paul's Notes}

\emph{Calculus} by Thomas \& Finney, 9th edition, pp 569-573

\end{document}
